% (C) 2026 Ralf Stephan, in collaboration with Claude Code.
% Released under CC0 1.0 Universal (public-domain dedication).
% See https://creativecommons.org/publicdomain/zero/1.0/

\documentclass[11pt]{article}
\usepackage[T1]{fontenc}
\usepackage[utf8]{inputenc}
\usepackage{lmodern}
\usepackage[margin=2.5cm]{geometry}
\usepackage{amsmath,amssymb}
\usepackage{braket}
\usepackage[hidelinks]{hyperref}

% Lean names and paths: verbatim, breakable at . _ /
\DeclareUrlCommand\lean{\urlstyle{tt}}
\def\UrlBreaks{\do\.\do\_\do\/}
% Marks an open problem that blocks checking a main claim of the brief
\newcommand{\blocker}[1]{\textsc{[blocks verification: #1]}~}
% Extra stretch for paragraphs with long, unbreakable Lean names
\emergencystretch=6em

\title{Open problems:\\ Quantum pure states are not enough}
\author{}
\date{}

\begin{document}
\maketitle

\noindent
Open problems raised by the technical brief in this directory, and by its Lean formalization
(see \lean{formalization-log.tex}).
Items marked \textsc{[blocks verification]} stand in the way of checking a main
claim of the brief as the brief states it: a \emph{formalization gap} can be closed
in Lean alone, while a \emph{brief issue} needs the brief to be corrected or made
precise first. Items marked \textsc{[optional]} would strengthen a result that is already
checked; no claim of the brief waits on them.

Theorem~5 is verified, and Theorem~3 for $\dim\mathcal H \ge 2$. Example~2 is verified at the
level of topology and measure and, on the Riemann sphere, of the symplectic form and of
Hamiltonian evolution, with the correction of the item ``Hamiltonian evolution on the classical
qubit''. Of Theorem~1 only the ``if'' for $V = \mathbb C^{n+1}$ is checked: see the two items
marked below.

\section*{From the brief}
\begin{enumerate}
\item \blocker{brief issue}\textbf{The range of Theorem~1.} The theorem claims consistency for
  $1 \le \dim V < \infty$, while its proof uses $2 \le \dim X$. Making it precise needs four
  choices.
  \begin{itemize}
  \item \emph{Dimension $0$.} The range starts at $1$ if a $0$-dimensional symplectic manifold
    ($n = 0$ degrees of freedom) is allowed, and at $2$ if not. For $\dim V = 1$ the state space
    is a single point, and $P(V)$ has two points exactly when $2 \le \dim V$
    (\lean{subsingleton_projectivization}, \lean{nontrivial_projectivization_iff}).
  \item \emph{Theorem~3 depends on it.} If a point is allowed, Theorem~3 is false for
    $\dim\mathcal H = 1$: the CST ensembles of a point and the density operators of
    $\mathbb C$ are both a single point, and the constant map identifies them
    (\lean{cst_qeDens_consistent_of_finrank_one}). So the brief cannot keep both
    $1 \le \dim V$ in Theorem~1 and Theorem~3 for every $\mathcal H$.
  \item \emph{The empty space.} For $\dim V = 0$, $P(V)$ is empty
    (\lean{isEmpty_projectivization}), which is a symplectic manifold of every dimension. Either
    lower bound, $1$ or $2$, needs CST to exclude it, for instance by requiring an ensemble.
  \item \emph{What ``represented by'' keeps.} If $X$ and $P(V)$ need only be in bijection, the
    ``only if'' is false: $P(\ell^2)$ and the sphere $P(\mathbb C^2)$ both have $\mathfrak c$
    points (\lean{nonempty_projectivization_equiv}), and the sphere's symplectic structure can be
    carried over. The theorem needs $V$ Hausdorff and $X$ to carry the quotient topology of
    $P(V)$; finiteness then comes from CST's Liouville measure $\omega^n$ with finite $n$, not
    from $X$ being symplectic (Log, ``Problems found'').
  \end{itemize}
  Proposed: exclude dimension $0$ and the empty space, so that Theorem~1 reads
  $2 \le \dim V < \infty$ for a Hausdorff $V$ with $X = P(V)$ as topological spaces. Theorem~3 and its Lean proof then stand
  as they are. The alternative keeps $1 \le \dim V$, also excludes the empty space and restricts
  Theorem~3 to $\dim\mathcal H \ge 2$. (Log, ``Problems found'', first three items.)
\item \textbf{Ensembles supported on a subspace.} The proof of Theorem~3 compares the ensembles
  with support in a two-dimensional subspace $A$. For $\dim\mathcal H \ge 3$ no CST ensemble has
  support in $P(A)$, which has Liouville measure zero. Whether the brief means the classical
  ensembles of the subsystem $P(A)$, with its own Liouville measure, is open. The formalization
  compares the whole ensemble spaces instead. (Log, ``Problems found'', second item.)
\item \textbf{Hamiltonian evolution on the classical qubit.} Example~2 identifies Hamiltonian
  evolution with the diffeomorphisms that conserve the volume. Area-preserving reflections of the
  sphere are not Hamiltonian, so the area form should be meant. Lean checks this for complex
  conjugation of the Riemann sphere, the Bloch reflection $a_1 \mapsto -a_1$
  (\lean{conjSphere_preserves_area_reverses_form}, \lean{globalFlow_ne_conjSphere},
  \lean{timeDependentFlowOf_ne_conjSphere}). With the area form the claim is checked for
  evolutions. A vector field preserves the Fubini--Study form exactly when it is Hamiltonian
  (\lean{lieDeriv_areaForm_eq_zero_iff_hamiltonian}), and a jointly smooth time-dependent flow
  preserves it exactly when every velocity field is Hamiltonian
  (\lean{preserves_areaForm_iff_hamiltonian}), as does a smooth isotopy from the identity through
  its generating field (\lean{isotopy_preserves_areaForm_iff_hamiltonian}). No flux homomorphism
  is needed: the obstruction is a class in $H^1_{dR}(S^2)$, which vanishes because the sphere is
  simply connected. Open is the
  statement for single maps, that every orientation-preserving, area-preserving diffeomorphism of
  the sphere is the end of a Hamiltonian evolution. It needs Smale's theorem
  $\mathrm{Diff}^+(S^2) \simeq SO(3)$, or at least that $\mathrm{Diff}^+(S^2)$ is connected through
  smooth isotopies, and a Moser argument for area forms on $S^2$. That quantum evolutions are
  exactly the rotations is checked (\lean{exists_unitaryAction_eq_iff}). (Log, ``Problems
  found'', Example~2 item.)
\end{enumerate}

\section*{From the formalization}
\begin{enumerate}
\item \blocker{formalization gap}\textbf{Theorem~1.} The ``if'' is formalized for
  $V = \mathbb C^{n+1}$: $\mathbb CP^n$, with TauCeti's real manifold structure, has exactly one
  symplectic form whose pullback along the Hopf map $S^{2n+1} \to \mathbb CP^n$ is the standard
  form $\operatorname{Im}\langle u, w\rangle$ of $\mathbb C^{n+1}$ restricted to the unit sphere,
  the Fubini--Study form (\lean{existsUnique_isSymplectic_projectivization}). Three parts are open.
  \begin{itemize}
  \item \emph{Arbitrary $V$.} TauCeti puts a manifold structure only on
    $P(\mathbb C^{n+1})$. For a complex vector space $V$ of finite dimension with a Hermitian
    inner product, the structure and the form are carried over along an orthonormal basis
    $V \simeq \mathbb C^{n+1}$. Two bases differ by a unitary map, which preserves the standard
    form and commutes with the Hopf map, so the form does not depend on the choice; this is not
    yet checked.
  \item \emph{The Liouville measure.} CST asks for the measure $\omega^n$ of the symplectic form,
    which is still missing upstream.
  \item \emph{The ``only if''.} That a symplectic $P(V)$ forces $\dim V < \infty$ needs the
    topology of $P(V)$; as sets it is false (\lean{nonempty_projectivization_equiv}). Theorem~3
    for $\dim\mathcal H = \infty$ rests on it, so it is formalized only for
    $\dim\mathcal H < \infty$.
  \end{itemize}
  (Log, ``Results''.)
\item \textbf{Example~2.} $P(\mathbb C^2)$ is homeomorphic to the sphere
  (\lean{nonempty_homeomorph_sphere}), its area is $dz\,d\theta$ (\lean{toSphere_eq_dz_dtheta}),
  the quantum evolutions are the rotations (\lean{exists_unitaryAction_eq_iff}) and every
  Hamiltonian makes the Bloch vector precess (\lean{exists_blochVector_precession}). Still
  missing: the area as the Liouville measure of the Fubini--Study form, and
  $\omega = dz \wedge d\theta$ in the chart, which need a Liouville measure upstream; the
  diffeomorphism $\mathbb CP^1 \cong S^2$ that compares the Fubini--Study form with the round area
  form, and TauCeti's form on the Riemann sphere, of total area $1$, with the form of the previous
  item for $n = 1$ (\lean{Projectivization.fubiniStudyForm}), which should give area $\pi$; the integrated rotation of the precession; the classical side of linear Hamiltonians,
  $X_{n\cdot x} = n \times x$, which needs moment maps. Refinements of Hamiltonian evolution are
  in the next item. (Log, ``Differences''.)
\item \textsc{[optional]}~\textbf{Refinements of Hamiltonian evolution.} The evolutions of
  Example~2 are checked with each velocity field $X_t$ Hamiltonian for its own $H_t$
  (\lean{preserves_areaForm_iff_hamiltonian}, \lean{isotopy_preserves_areaForm_iff_hamiltonian}).
  Three refinements remain.
  \begin{itemize}
  \item \emph{One jointly smooth $H(t, x)$.} TauCeti builds the primitive of a closed $1$-form
    by lifting through a covering space, separately for each $t$, and has no version with
    parameters. A route: normalize $H_t(x_0) = 0$; prove a Poincaré lemma with a time parameter
    in charts, where the flat step is the homotopy operator \lean{poincareHomotopy} made jointly
    smooth with \lean{contDiffOn_integral_smul_of_contDiffOn}; then the points $y$ at which
    $t \mapsto H_t(y)$ is smooth form an open and closed set, and $H$ is jointly smooth near
    each point. The chart step repeats TauCeti's manifold Poincaré lemma with a parameter and
    needs chart expressions of $\iota_{X_t}\omega$; estimated at 400--700 lines.
  \item \emph{Area instead of area form along an evolution.} Along an isotopy from the identity,
    $\varphi_t^*\omega = \pm\omega$ forces the sign $+$, since it cannot jump; with area
    preservation stated pointwise this is a short continuity argument. The measure-theoretic
    statement needs the Liouville measure of $\omega$, which is missing upstream (item~2).
  \item \emph{Conservation of entropy.} The brief's ``and therefore the entropy is conserved''
    follows from measure preservation, since the Gibbs entropy is invariant under
    measure-preserving bijections. For the rotations this is direct; for a general Hamiltonian
    flow it needs the Liouville measure, as above.
  \end{itemize}
\item \textbf{Extreme points.} The proof of Theorem~5 starts from the fact that the pure states are
  the extreme points of $D(\mathcal H)$. The formalization does not need it and does not prove
  it. (Log, ``Differences'', last item.)
\end{enumerate}

\end{document}
